\subsection{伪同构}

保留第 2、3 节的记号和假设。

命题 9. 设 M 是有限生成的 A-模。下列条件等价：

\begin{enumerate}
\def\labelenumi{(\alph{enumi})}
\item
  对每个高度  的素理想 ，有 \(\mathbf{M}_{\mathfrak{p}} = 0\)。\(\mathfrak{p}\)\(\leqslant 1\)
\item
   的湮灭子 \(\mathfrak{a}\) 是非零理想，且 \(\mathbf{M}\)\(\neq (0)\)\(\mathbf{A}:\mathfrak{a} = \mathbf{A}\)（如 § 1，第 1 节，\(\mathbf{A}:\mathfrak{a}\) 表示由满足  的 \(x\in \mathbf{K}\) 组成的集合）。\(xa\in \mathbf{A}\)
\end{enumerate}

我们知道（第二章，§ 2，第 2 节，命题 4 的推论 2），条件 \(\mathbf{M}_{\mathfrak{p}} = 0\) 等价于 \(\mathfrak{a} \nsubseteq \mathfrak{p}\)，因而等价于 \(\mathfrak{a}\mathbf{A}_{\mathfrak{p}} = \mathbf{A}_{\mathfrak{p}}\)（第二章，§ 2，第 5 节，注）；另一方面，对于 A 的每个非零整理想 \(\mathfrak{b} \neq 0\)，关系``对所有 ，\(\mathfrak{b}\mathbf{A}_{\mathfrak{p}} = \mathbf{A}_{\mathfrak{p}}\)''等价于 D(A) 中的 \(\mathfrak{p} \in \mathbb{P}\)\(\operatorname{div} \mathfrak{b} = \operatorname{div} \mathbf{A} = 0\)（§ 1，第 4 节，命题 7），或者也等价于 \(\operatorname{div}(\mathbf{A}: \mathfrak{b}) = 0\)；由于 A: \(\mathfrak{b}\) 是除性理想（§ 1，第 1 节，命题 1），这个关系还等价于 A: \(\mathfrak{b} = \mathbf{A}\)。注意到当  时说 \(\mathfrak{a} \nsubseteq \mathfrak{p}\) 意味着 \(\mathfrak{p} = (0)\)\(\mathfrak{a} \neq (0)\)，命题得证。

注（1）。命题 9 的等价条件还意味着 \(\operatorname{Ass}(\mathbf{M})\) 不包含高度 \(\leqslant 1\) 的素理想。* 也可以将其解释为：\(\operatorname{Supp}(\mathbf{M})\) 在  中的余维 \(\geqslant 2\)。*\(\operatorname{Spec}(\mathbf{A})\)

定义 2. 若 A-模 M 是有限生成的且满足命题 9 的等价条件，则称其为伪零模。

该定义和命题 9 表明，伪零 A-模是挠 A-模；反之不成立。

\textbf{例子}\par

\begin{enumerate}
\def\labelenumi{(\arabic{enumi})}
\item
  若 A 是 Dedekind 整环，则 A 的每个素理想的高度都 \(\leqslant1\)；因此说 M 是伪零的就意味着 Supp(M) = \(\varnothing\)，从而 M = 0（第二章，§ 4，第 4 节）。
\item
  令 \(k\) 为一个域，\(A = k[X, Y]\) 为 \(k\) 上两个不定元的多项式环；若 \(m\) 是  的极大理想 \(AX + AY\)，则 \(A\)\(A\)-模 \(A/m\) 是伪零的；因为其湮灭子 \(m\) 的高度不为 \(\leqslant 1\)，这是由于它包含主素理想 \(AX\) 和 \(AY\)，且不同于二者；因此 \(A: m = A\)（§ 1，第 6 节，定理 3 的推论 1）。
\end{enumerate}

定义 3. 设 M 和 N 是两个 A-模，f: M → N 是一个同态。若 Ker(f)（相应地 Coker(f)、Im(f)）是伪零的，则称 f 为伪单射（相应地为伪满射、伪零）；若 f 既是伪单射又是伪满射，则称 f 为伪双射。

伪双射同态也称为伪同构。

设 M 和 N 是有限生成的；则 \(f: M \to N\) 为伪单射（相应地为伪满射、伪零）的充要条件是，对所有 \(p \in P \cup \{\{0\}\}, f_{p}: M_{p} \to N_{p}\)， 是单射（相应地是满射、零映射）；这由 A-模 \(A_{p}\) 的平坦性得出（参见第一章，§ 2，第 3 节，注 2）。

例（3）。设 M 是有限生成的无挠 A-模；则 M 到其二重对偶的典范映射 \(c_{M}: M \to M^{**}\) 是伪同构。因为 M 可识别为 \(V = M \otimes_{A} K\) 的格（第 1 节，命题 1）；我们已经看到，对所有 ，\(M_{p} = M_{p}^{**}\)（第 2 节，例 2），而当 p = 0 时，\(p \in P\)\(M_{p}\) 和 \(M_{p}^{**}\) 都等于 V。

定理 4. 设 E 是有限生成的 A-模，T 是 E 的挠子模，且 M = E/T。存在一个伪同构

\[
f \colon \mathbf {E} \rightarrow \mathbf {T} \times \mathbf {M}.
\]

我们先证明两个引理。

引理 2. 设 \((\mathfrak{p}_i)_{1 \leqslant i \leqslant k}\) 是 A 的高度为 1 的素理想组成的非空有限族，并令 \(S = \bigcap_{i} (A - \mathfrak{p}_i)\)；则环 \(S^{-1}A\) 是主理想整环。

\(S^{-1}A\) 是半局部环，其极大理想为 \(m_{i} = p_{i}S^{-1}A\)（\(1 \leqslant i \leqslant k\)），局部环 \((S^{-1}A)_{m_{i}}\) 同构于 \(A_{p_{i}}\)（第二章，§3，第 5 节，命题 17），因而是离散赋值环。因此，环 \(S^{-1}A\) 是 Dedekind 整环（§2，第 2 节，定理 1 (f)；又由于它是半局部环，故是主理想整环（§2，第 2 节，命题 1）。

引理 3. 存在一个同态 \(g: \mathbf{E} \to \mathbf{T}\)，其在 \(\mathbf{T}\) 上的限制既是伸缩同态又是伪同构。

令 \(\mathfrak{a}\) 为 \(\mathrm{T}\) 的湮灭子；由于 \(\mathrm{T}\) 是有限生成的挠 A-模，\(\mathfrak{a} \neq 0\)。令 \(\mathfrak{p}_i\)（\(1 \leqslant i \leqslant k\)）为包含 \(\mathfrak{a}\) 的高度为 1 的素理想（其数目有限（§ 1，第 6 节，定理 4））；若此数为 0，则 T 是伪零的（命题 9(a)），可以取 g = 0。否则，令 \(S = \bigcap_{i} (A - p_{i})\)；由引理 2，\(S^{-1}A\) 是主理想整环，因而 \(S^{-1}M\)（这是无挠的有限生成 \(S^{-1}A\)-模）是自由的（代数，第七章，§ 4，第 3 节，定理 2 的推论 2），又由于 \(S^{-1}M = (S^{-1}E)/(S^{-1}T)\)，\(S^{-1}T\) 是 \(S^{-1}E\) 的直因子（代数，第二章，§ 1，第 11 节，命题 21）。现在，

\[
\operatorname{Hom} _ {\mathbf {S} ^ {- 1} \mathbf {A}} (\mathrm{S} ^ {- 1} \mathrm{E}, \mathrm{S} ^ {- 1} \mathrm{T}) = \mathrm{S} ^ {- 1} \operatorname{Hom} _ {\mathbf {A}} (\mathrm{E}, \mathrm{T})
\]

（第二章，§2，第 7 节，命题 19）；因此存在 \(s_0 \in S\) 和 \(g_0 \in \mathrm{Hom}_{\mathbf{A}}(\mathbf{E}, \mathbf{T})\)，使得 \(s_0^{-1}g_0\) 是从 \(\mathbf{S}^{-1}\mathbf{E}\) 到 \(\mathbf{S}^{-1}\mathbf{T}\) 的投影算子。若 \(h_0 \in \mathrm{Hom}_{\mathbf{A}}(\mathbf{T}, \mathbf{T})\) 表示 \(g_0\) 在 \(\mathbf{T}\) 上的限制，则存在 \(s_1 \in \mathbf{S}\)，使得对所有  有 \(s_1h_0(x) = s_1s_0x\)；令 \(x \in \mathbf{T}\)\(s = s_1s_0\)、\(g = s_1g_0\)、\(h = s_1h_0\)，则 h 是 T 上比为 s 的伸缩，并且是 g 在 T 上的限制。还需验证 h 是伪同构。现在，若 \(h\)\(s\)\(\mathbf{T}\)\(g\)\(\mathbf{T}\)\(h\)\(p = 0\)，或 \(p \in P\) 不同于 \(p_i (1 \leqslant i \leqslant k)\)，则 \(T_p = 0\)（第二章，§4，第 4 节，命题 17），且 \(h_p: T_p \to T_p\) 是同构；反之，若 \(p\) 等于某个 \(p_i (1 \leqslant i \leqslant k)\)，则 s 在 \(s\)\(A_{p_i}\) 中可逆，而 \(h_{p_i}\)（T＿\(s\)\(T_{p_i}\)p＿i 上比为 s 的伸缩）也是同构，这完成了引理 3 的证明。

现在证明定理 4。令 \(g: \mathbf{E} \to \mathbf{T}\) 为满足引理 3 性质的同态；令 h 为 g 在 \(h\)\(g\)\(\mathbf{T}\) 上的限制，令 \(\pi\) 为 \(\mathbf{E}\) 到 \(\mathbf{M}\) 的典范投影。我们证明同态 \(f = (g, \pi): \mathbf{E} \to \mathbf{T} \times \mathbf{M}\) 解决了问题。交换图为：

\[
\begin{array}{c c c} 0 \longrightarrow \mathrm{T} \longrightarrow & \mathrm{E} & \longrightarrow \mathrm{M} \longrightarrow 0 \\ h \Big \downarrow & f \Big \downarrow & ^ {1 _ {\mathrm{M}}} \Big \downarrow \\ 0 \longrightarrow \mathrm{T} \longrightarrow \mathrm{T} \times \mathrm{M} \longrightarrow \mathrm{M} \longrightarrow 0 \end{array}
\]

其中两行都是正合的。蛇形图（第一章，§ 1，第 4 节，命题 2）给出正合列：

\[
0 \rightarrow \operatorname{Ker} (h) \rightarrow \operatorname{Ker} (f) \rightarrow 0 \rightarrow \operatorname{Coker} (h) \rightarrow \operatorname{Coker} (f) \rightarrow 0
\]

从而 \(\operatorname{Ker}(f)\) 同构于 \(\operatorname{Ker}(h)\)，而 \(\operatorname{Coker}(f)\) 同构于 \(\operatorname{Coker}(h)\)。由于 h 是伪同构，f 也是伪同构。\(h\)\(f\)

可以说，在``伪同构意义下''，定理 4 将有限生成 A-模的研究归结为两部分：一方面是无挠模的研究，另一方面是挠模的研究。此外，我们在上面（例 3）已经看到，无挠模与其二重对偶伪同构，因而与一个自反模伪同构。至于挠模，有如下结果，它们在伪同构意义下被确定：

定理 5. 设 T 是有限生成的挠 A-模。存在两个有限族 \((n_i)_{i\in I}\) 和 \((p_i)_{i\in I}\)，其中 \(n_i\) 是满足 \(\geqslant 1\) 的整数，\(p_i\) 是 A 的高度为 1 的素理想，并且若写成 \(\mathrm{T}' = \bigoplus_{i\in I} \mathrm{A} / p_i^{n_i}\)，则存在从 T 到 \(\mathrm{T}'\) 的伪同构。此外，具有此性质的族 \((n_i)_{i\in I}\) 和 \((p_i)_{i\in I}\) 在指标集的一个双射意义下是唯一的，且 \(p_i\) 包含 T 的湮灭子。

唯一性：若 \(f\colon T\to T'\) 是伪同构，且对 \(\mathfrak{p}\in\mathbf{P},f_{\mathfrak{p}}\colon T_{\mathfrak{p}}\to T'_{\mathfrak{p}}\)， 是同构。现在，\(T'_{\mathfrak{p}}\) 是各个 \(A_{\mathfrak{p}}/p^{n_{i}}A_{\mathfrak{p}}\) 的直和，求和遍历满足  的指标 \(i\)；因此 \(\mathfrak{p}_{i}=\mathfrak{p}\)\(p^{n_{i}}A_{\mathfrak{p}}\) 是挠 \(A_{\mathfrak{p}}\)-模 \(T_{\mathfrak{p}}\) 的初等因子（代数，第七章，§ 4，第 7 节），其唯一性已在代数，第七章，§ 4，第 7 节，命题 7 中证明。

存在性：可将注意力限于 \(\mathbf{T} \neq 0\) 的情形。令 \(\mathfrak{a}\) 为 \(\mathbf{T}\) 的湮灭子（非零且不同于 A），令 \(\mathfrak{p}_i\)（\(1 \leqslant i \leqslant k\)）为包含 \(\mathfrak{a}\) 的 A 的高度为 1 的素理想（其数目有限（§ 1，第 6 节，定理 4）），并令 \(\mathbf{S} = \bigcap_i (\mathbf{A} - \mathfrak{p}_i)\)。半局部环 \(\mathbf{A}' = \mathbf{S}^{-1}\mathbf{A}\) 是主理想整环（引理 2），其极大理想为 \(\mathfrak{m}_i = \mathfrak{p}_i\mathbf{A}'\)；由于 \(\mathbf{S}^{-1}\mathbf{T}\) 是有限生成的挠 \(\mathbf{A}'\)-模，它同构于有限直和 \(\bigoplus_{f \in \mathbf{I}} \mathbf{A}' / \mathfrak{m}_{\phi(f)}^{n_f}\)，其中 \(\phi\) 是从有限集 I 到 \([1, k]\) 的映射（代数，第七章，§ 4，第 7 节，命题 7）；由于 \(\mathbf{A}' / \mathfrak{m}_{\phi(f)}^{n_f}\) 同构于 \(\mathbf{S}^{-1}(\mathbf{A} / \mathfrak{p}_{\phi(f)}^{n_f})\)（第二章，§ 2，第 4 节），我们得到所需类型的挠 A-模 \(\mathbf{T}'\)，以及把 \(f_0\)\(\mathbf{S}^{-1}\mathbf{T}\) 映到 \(\mathbf{S}^{-1}\mathbf{T}'\) 上的同构 。由于 \(\operatorname{Hom}_{\mathbf{S}^{-1}\mathbf{A}}(\mathbf{S}^{-1}\mathbf{T}, \mathbf{S}^{-1}\mathbf{T}')\) 等于 \(\mathbf{S}^{-1}\operatorname{Hom}_{\mathbf{A}}(\mathbf{T}, \mathbf{T}')\)（第二章，§ 2，第 7 节，命题 19），存在 \(s \in S\) 和同态 \(f: T \to T'\)，使得 \(f_0 = s^{-1}f\)。还需证明 f 是伪同构：现在，若 \(f\)\(\mathfrak{p} = 0\)，或 \(\mathfrak{p} \in P\) 不同于 \(\mathfrak{p}_i\)，则 \(\mathrm{T}_{\mathfrak{p}} = \mathrm{T}_{\mathfrak{p}}' = 0\)（第二章，§ 4，第 4 节，命题 17）；反之，若 \(\mathfrak{p}\) 是某个 \(\mathfrak{p}_i\)（\(1 \leqslant i \leqslant k\)），则 s 在 \(s\)\(\mathbf{A}_{\mathfrak{p}_i}\) 中可逆，并且由于 \(f_{\mathfrak{p}_i} = s(f_0)_{\mathfrak{p}_i}\)，而 \((f_0)_{\mathfrak{p}_i}\) 是从 \(\mathrm{T}_{\mathfrak{p}_i} = (\mathrm{S}^{-1}\mathrm{T})_{\mathfrak{m}_i}\)\(\mathrm{T}_{\mathfrak{p}_i}' = (\mathrm{S}^{-1}\mathrm{T}')_{\mathfrak{m}_i}\) 到  的同构，所以 \(f_{\mathfrak{p}_i}\) 也是同构。

注（2）。在定理 5 的陈述中，模 \(\mathbf{A} / \mathfrak{p}_i^n\) 可以替换为 \(\mathbf{A} / \mathfrak{p}_i^{(n_i)}\)（§ 1，第 4 节，命题 8）。对于所有 \(\mathfrak{p} \in \mathbb{P}\)，典范映射 \(g\colon \mathbf{A} / \mathfrak{p}^n \to \mathbf{A} / \mathfrak{p}^{(n)} = \mathbf{A} / (\mathbf{A} \cap \mathfrak{p}^n \mathbf{A}_{\mathfrak{p}})\) 是伪同构，因为对于所有不同于  的 \(\mathfrak{q} \in \mathbb{P}\)，\(\mathfrak{p}\)\(\mathbf{A}_{\mathfrak{q}} / \mathfrak{p}^n \mathbf{A}_{\mathfrak{q}} = \mathbf{A}_{\mathfrak{q}} / \mathfrak{p}^{(n)} \mathbf{A}_{\mathfrak{q}} = 0\)，且 \(\mathbf{A}_{\mathfrak{p}} / \mathfrak{p}^n \mathbf{A}_{\mathfrak{p}} = \mathbf{A}_{\mathfrak{p}} / \mathfrak{p}^{(n)} \mathbf{A}_{\mathfrak{p}}\)。

* 给定 A-模的正合列 E → F → G，若 E 和 G 是伪零的，则 F 也是伪零的，这由定义 2 和第二章，§ 2，第 4 节，定理 1 得出。用范畴语言，我们于是可以说，在 A-模范畴 C 中，伪零模的子范畴 C' 是满子范畴，从而可以定义商范畴 C/C'：该范畴的对象仍是 A-模，但从 E 到 F 的态射集合（其中 E、F 属于 C/C'）是集合 \(\mathrm{Hom}_{\mathbf{A}}(\mathbf{E}',\mathbf{F}')\) 的直极限，其中 \(\mathbf{E}'\)（相应地 \(\mathbf{F}'\)）分别遍历 \(\mathbf{E}\) 的子模集合（相应地， 的商模 \(\mathbf{F} / \mathbf{F}''\) 的集合），使得 \(\mathbf{F}\)\(\mathbf{E} / \mathbf{E}'\)（相应地 \(\mathbf{F}''\)）是伪零的。当然，对于每一对 A-模 \(\mathbf{E}\)、\(\mathbf{F}\)，都有典范同态 \(\mathrm{Hom}_{\mathcal{C}}(\mathbf{E},\mathbf{F})\to \mathrm{Hom}_{\mathcal{C} / \mathcal{C}'}(\mathbf{E},\mathbf{F})\)。称同态 \(u\in \mathrm{Hom}_{\mathbf{A}}(\mathbf{E},\mathbf{F})\) 是伪零的（相应地是伪单射、伪满射、伪双射），是指它在 \(\mathrm{Hom}_{\mathcal{C} / \mathcal{C}'}(\mathbf{E},\mathbf{F})\) 中的典范像为零（相应地为单态、满态、同构）。*

